\chapter{La coquille d'application : ce que vous n'écrivez plus}

Le chapitre précédent vous a fait écrire une application \emph{à la main} :
créer la fenêtre, ouvrir le contexte graphique, tenir la boucle, écouter les
événements, présenter l'image. C'était la bonne façon de comprendre ce qui se
passe --- et c'est la seule façon d'apprendre ce qu'une coquille vous cache.

Ce chapitre montre la coquille. Depuis \texttt{NkCanvasApp}, cette plomberie
devient quelques méthodes à remplir. Surtout, la coquille apporte trois services
qu'une application écrite à la main \emph{oublie presque toujours} --- et l'oubli
ne se voit qu'une fois le programme entre les mains de quelqu'un d'autre.

\begin{nkretenir}
Une coquille ne vous fait pas gagner des lignes : elle vous fait gagner les
lignes que \emph{vous n'auriez pas écrites}.

Mesure du dépôt, le 2 septembre 2026 : \textbf{quatorze} applications créent
elles-mêmes leur fenêtre de rendu. Sur ces quatorze, \textbf{trois} gèrent le
cycle de vie mobile. Les onze autres --- dont \texttt{Sandbox}, \texttt{NkRef},
\texttt{NKGuiDemo}, \texttt{NkVideoPlayer} et \texttt{NkAudioPlayer} --- se
ferment quand le téléphone se verrouille, et personne ne l'a remarqué parce que
personne ne les lance sur un téléphone.

Ce n'est pas de la négligence : c'est que \emph{rien ne rappelle d'écrire ce
qu'on ne voit pas manquer}. C'est précisément ce qu'une coquille existe pour
donner.
\end{nkretenir}

\section{Le plus petit programme complet}

\begin{nkcode}{Applications/NkDames/src/main.cpp}
#include "Dames/NkDamesJeu.h"
#include "NKWindow/Core/NkEntry.h"
#include "NKWindow/NKMain.h"

NKENTSEU_DEFINE_APP_DATA(([]() {
    nkentseu::NkAppData d{};
    d.appName = "NkDames";
    d.appVersion = "1.0.0";
    return d;
})());

int nkmain(const nkentseu::NkEntryState &state) {
    return nkentseu::renderer::NkCanvasApp::Run<nkentseu::jeux::dames::NkDamesJeu>(state);
}
\end{nkcode}

Trois choses seulement, et elles se lisent dans l'ordre.

\texttt{NKENTSEU\_DEFINE\_APP\_DATA} déclare le nom et la version de
l'application. \texttt{nkmain} est le point d'entrée portable : c'est
\texttt{NKWindow/NKMain.h} qui fournit derrière lui le \texttt{WinMain} de
Windows et l'\texttt{android\_main} d'Android. \texttt{NkCanvasApp::Run<T>}
construit votre classe et lui rend la main.

\begin{nkpiege}
La macro attend un \texttt{NkAppData}, \textbf{pas une chaîne} --- son nom laisse
croire le contraire. Écrivez \texttt{NKENTSEU\_DEFINE\_APP\_DATA("MonJeu")} et le
compilateur répond \texttt{no viable overloaded '='}, sans jamais nommer la
cause. Le lambda ci-dessus est le patron employé par toutes les applications du
dépôt.
\end{nkpiege}

\section{Les méthodes que vous remplissez}

\begin{nkcode}{Kernel/Runtime/NKCanvas/src/NKCanvas/App/NkCanvasApp.h (extraits)}
virtual NkOptional<int> OnCommandLine(const NkVector<NkString> &args);
virtual bool OnInit();
virtual void OnUpdate(float32 deltaTime);
virtual void OnRender(NkRenderWindow &target);
virtual bool OnEvent(const NkEvent &event);      // true = consomme
virtual bool OnPointer(const NkPointer &pointer);
virtual void OnLayout(const NkLayoutInfo &layout);
virtual bool OnCloseRequested();                 // false = refuse la fermeture
virtual void OnPause();
virtual void OnResume();
virtual void OnShutdown();
\end{nkcode}

Aucune n'est obligatoire : elles ont toutes un corps par défaut. On remplit ce
dont on a besoin.

\begin{center}
\begin{tabular}{@{}ll@{}}
\toprule
\textbf{Méthode} & \textbf{Quand elle est appelée} \\
\midrule
\texttt{OnCommandLine} & avant toute fenêtre --- peut refuser de démarrer \\
\texttt{OnInit}        & la fenêtre et le GPU existent \\
\texttt{OnUpdate}      & une fois par image, avec le pas de temps \\
\texttt{OnRender}      & une fois par image, avec la cible de rendu \\
\texttt{OnPointer}     & souris \emph{ou} doigt, unifiés \\
\texttt{OnLayout}      & taille, densité et zone sûre ont changé \\
\texttt{OnPause}, \texttt{OnResume} & l'application part en arrière-plan, revient \\
\bottomrule
\end{tabular}
\end{center}

\begin{nkpiege}
\texttt{OnCommandLine} rend un \texttt{NkOptional<int>}. Rendre une valeur
\emph{arrête l'application avec ce code}, sans ouvrir de fenêtre --- c'est ce qui
permet à une option de banc de répondre en mode automatique. Rendre un
\texttt{NkOptional} vide laisse le programme démarrer.
\end{nkpiege}

\section{Les trois services que la coquille rend}

\subsection{La boucle cède la main, et sur le Web c'est vital}

\begin{nkcode}{Kernel/Runtime/NKCanvas/src/NKCanvas/App/NkCanvasApp.cpp (CadencerTrame)}
if (mConfig.imagesParSeconde <= 0) {
    NkChrono::YieldThread();
    return;
}
const float64 budgetMs = 1000.0 / static_cast<float64>(mConfig.imagesParSeconde);
const float64 ecouleMs = chrono.Elapsed().milliseconds;
if (ecouleMs < budgetMs) {
    NkChrono::Sleep(budgetMs - ecouleMs);
} else {
    NkChrono::YieldThread();
}
\end{nkcode}

Sous Emscripten, \texttt{NkChrono::Sleep} appelle \texttt{emscripten\_sleep(ms)}
et \texttt{YieldThread} appelle \texttt{emscripten\_sleep(0)}. Les deux rendent
la main au navigateur.

\begin{nkpiege}
\textbf{Les deux branches cèdent, et c'est tout l'intérêt.} Une version
« ne dormir que si on est en avance » gèle l'onglet dès la première image lente
--- c'est-à-dire tout le temps sur le Web, où une image dépasse presque toujours
son budget. Le \texttt{else} n'est pas une optimisation : c'est lui qui rend la
main.

Le mécanisme existait avant la coquille, et il était juste. Ce qui manquait,
c'était l'\emph{appel}.
\end{nkpiege}

\subsection{La zone sûre, mesurée et non devinée}

\begin{nkcode}{Kernel/Runtime/NKCanvas/src/NKCanvas/App/NkCanvasApp.h (NkLayoutInfo)}
struct NkLayoutInfo {
    uint32 width  = 0;
    uint32 height = 0;
    float32 density = 1.f;   // facteur d'echelle de l'ecran (DPI)
    NkSafeAreaInsets safeArea;
};
\end{nkcode}

\texttt{safeArea} porte quatre marges --- \texttt{top}, \texttt{bottom},
\texttt{left}, \texttt{right} --- en pixels. Elles décrivent la partie de l'écran
qui n'est masquée ni par l'encoche, ni par les coins arrondis, ni par
l'indicateur de geste du bas.

\begin{nkretenir}
La zone sûre est une \textbf{ancre}, pas une marge : elle change à la rotation,
à l'ouverture du clavier, en écran partagé, et d'un appareil à l'autre. On la
demande à la plateforme à l'exécution ; une marge écrite en dur est fausse sur le
modèle suivant.

Et c'est un choix \emph{par élément} : les fonds, les images de couverture et les
listes qui défilent vont jusqu'au bord ; le texte et les boutons restent dans la
zone sûre. Un fond qui s'arrête à la zone sûre laisse des bandes disgracieuses ;
un bouton qui déborde sous l'indicateur de geste est \textbf{inatteignable}.
\end{nkretenir}

\subsection{Le pointeur : souris et doigt unifiés}

\begin{nkcode}{Kernel/Runtime/NKEvent/src/NKEvent/NkPointerEvent.h}
enum class NkPointerPhase : uint8 {
    NK_POINTER_NONE = 0,   // l'evenement lu ne concerne pas le pointeur
    NK_POINTER_DOWN,       // clic gauche enfonce, ou doigt pose
    NK_POINTER_MOVE,       // deplacement
    NK_POINTER_UP,         // clic relache, ou doigt leve
    NK_POINTER_CANCEL      // le systeme a repris la main (appel entrant, geste OS)
};

struct NkPointer {
    float32 x = 0.f;
    float32 y = 0.f;
    NkPointerPhase phase = NkPointerPhase::NK_POINTER_NONE;
    uint32 id = 0;            // 0 = souris, ou premier doigt
    bool fromTouch = false;   // true = doigt

    bool IsValid() const noexcept;   // phase != NK_POINTER_NONE
};
\end{nkcode}

Un clic de souris et un appui du doigt arrivent dans la même méthode, avec la
même forme. \texttt{fromTouch} dit lequel c'était, pour les rares cas où la
distinction compte --- une infobulle au survol n'a pas de sens au doigt.

\begin{nkpiege}
\textbf{\texttt{x} et \texttt{y} sont en pixels \emph{client}} : le repère de la
zone de dessin, origine en haut à gauche, hors bordure de fenêtre. C'est le même
repère qu'une cible de rendu. N'y mêlez jamais des coordonnées écran --- le dépôt
a payé une journée entière de recherche sur une sélection morte, pour avoir
comparé des pixels de fenêtre à des pixels de viseur.
\end{nkpiege}

\begin{nkpiege}
\textbf{Agissez au relâchement, pas à l'appui.} \texttt{NK\_POINTER\_UP} permet à
l'utilisateur de glisser \emph{hors} d'un bouton pour annuler son geste ---
comportement attendu partout, et gratuit ici. Un bouton qui déclenche à l'appui
ne pardonne pas l'erreur de visée, et sur un téléphone on vise mal.

Et n'oubliez pas \texttt{NK\_POINTER\_CANCEL} : un appel entrant ou un geste
système interrompt la saisie. Le traiter comme un relâchement laisse un objet
« collé » au doigt qui n'est plus là.
\end{nkpiege}

\section{Deux coquilles, et il faut choisir}

\texttt{NkCanvasGuiApp} dérive de \texttt{NkCanvasApp} et y ajoute NKGui : le
contexte, les polices, le backend de dessin. Mais elle \emph{prend la place} du
rendu.

\begin{nkcode}{Kernel/Runtime/NKCanvas/src/NKCanvas/App/NkCanvasGuiApp.h}
void OnUpdate(float32 deltaTime) final {
    mLastDelta = deltaTime;
    OnTick(deltaTime);
}

void OnRender(NkRenderWindow &target) final {
    const math::NkVec2u size = target.GetSize();
    mGuiContext.BeginFrame(mLastDelta);
    OnDraw(mGuiContext.dl);
    mGuiContext.EndFrame();
    mGuiBackend.Submit(mGuiContext.dl, size.x, size.y);
    mGuiBackend.Submit(mGuiContext.dlOverlay, size.x, size.y);
}
\end{nkcode}

Les deux méthodes sont \texttt{final}. Une classe qui dérive de
\texttt{NkCanvasGuiApp} ne peut donc \textbf{pas} redéfinir \texttt{OnRender} :
elle remplit \texttt{OnTick} et \texttt{OnDraw} à la place, et elle dessine avec
une \texttt{NkGuiDrawList}, pas avec un \texttt{NkRenderer2D}.

\begin{center}
\begin{tabular}{@{}lll@{}}
\toprule
\textbf{Vous dérivez de} & \textbf{Vous dessinez avec} & \textbf{Vous remplissez} \\
\midrule
rien (à la main)        & \texttt{NkRenderer2D}   & tout \\
\texttt{NkCanvasApp}    & \texttt{NkRenderer2D}   & \texttt{OnUpdate}, \texttt{OnRender} \\
\texttt{NkCanvasGuiApp} & \texttt{NkGuiDrawList}  & \texttt{OnTick}, \texttt{OnDraw} \\
\bottomrule
\end{tabular}
\end{center}

\begin{nkretenir}
\textbf{Comment choisir.} Vous dessinez des formes, des sprites, un monde de
jeu : \texttt{NkCanvasApp} et \texttt{NkRenderer2D}. Vous dessinez une interface
--- boutons, listes, panneaux : \texttt{NkCanvasGuiApp} et la liste
d'affichage. Vous voulez comprendre ce qui se passe dessous, ou vous avez un
besoin qu'aucune des deux ne couvre : à la main, comme au chapitre précédent.

Ce ne sont pas trois niveaux de qualité : ce sont trois outils pour trois
travaux.
\end{nkretenir}

\section{L'écran d'ouverture, offert}

\begin{nkcode}{Applications/NkLudo/src/Ludo/NkLudoJeu.cpp}
bool NkLudoJeu::OnGuiInit() {
    mSplash.PoserJeu("Ludo");
    Rejouer();
    return true;
}

void NkLudoJeu::OnTick(float32 deltaTime) {
    if (mSplash.Avancer(deltaTime)) {
        return;                       // l'ouverture tient l'ecran
    }
    /* ... le jeu ... */
}
\end{nkcode}

\texttt{NkCanvasSplash} peint deux volets --- la marque Rihen, puis le moteur et
le nom du jeu --- avec le logo officiel embarqué, rasterisé depuis un SVG de
514 octets. \texttt{PoserJeu} est le \emph{seul} réglage : les couleurs, les
durées et la mise en page sont les mêmes partout, délibérément.

\begin{nkpiege}
\texttt{Avancer} rend \texttt{true} tant que l'ouverture tient l'écran, et il
faut \textbf{s'arrêter là}. Sans ce retour, le jeu tournerait \emph{derrière}
l'écran de marque : une IA jouerait ses premiers coups pendant les quatre
secondes d'ouverture, et le joueur trouverait la partie déjà entamée en arrivant.

Et l'ouverture doit \textbf{manger l'appui} qui la saute. Sans cela, le même
appui passe le splash \emph{et} active le bouton qui se trouve dessous ---
l'utilisateur lance une partie qu'il n'a pas choisie.
\end{nkpiege}

\begin{nkpratique}
Reprenez l'application du chapitre précédent, écrite à la main, et portez-la sur
\texttt{NkCanvasApp}. Comptez les lignes avant et après.

Puis répondez à deux questions, en les \emph{mesurant} plutôt qu'en les
devinant : votre version manuelle gérait-elle \texttt{OnPause} ? Et cédait-elle
la main à chaque image ? Compilez-la pour le Web et ouvrez-la dans un
navigateur --- vous aurez la réponse en une seconde.
\end{nkpratique}

\begin{nkbilan}
La coquille tient la fenêtre, la boucle, le cycle de vie mobile, la zone sûre et
le pointeur unifié --- c'est-à-dire tout le chapitre précédent. Vous ne
remplissez que \texttt{OnGuiInit}, \texttt{OnTick} et le dessin.
\texttt{OnRender} est \texttt{final} : c'est la seule chose qu'elle vous
interdit, et c'est ce qui garantit que les quatre gestes restent dans le bon
ordre. Deux coquilles existent, et le choix se fait sur une seule question :
avez-vous besoin de NKGui, ou seulement de dessiner ?
\end{nkbilan}

\begin{nkquestions}
\begin{enumerate}
\item Citez trois services que la coquille rend et que vous n'écrivez plus.
\item Pourquoi \texttt{OnRender} est-il \texttt{final} ?
\item Qu'est-ce que la zone sûre, et pourquoi ne se code-t-elle pas en dur ?
\item Sur le Web, que se passe-t-il si la boucle ne cède pas la main ?
\item Quelle question départage \texttt{NkCanvasApp} et \texttt{NkCanvasGuiApp} ?
\end{enumerate}
\end{nkquestions}

\begin{nkdemo}
Prouvez que la zone sûre est \emph{mesurée} et non devinée.

Affichez, à chaque image, les quatre marges de \texttt{NkLayoutInfo} en clair sur
l'écran. Lancez sur ordinateur : elles doivent être nulles ou presque. Lancez sur
téléphone, puis faites pivoter l'appareil.

\textbf{Ce que la démonstration établit} : les valeurs \emph{changent} à la
rotation. Une marge codée en dur aurait été juste dans une orientation et fausse
dans l'autre --- et le défaut ne se voit que sur l'appareil, là où il coûte le
plus cher.
\end{nkdemo}

\begin{nklogique}
Votre écran d'ouverture dure quatre secondes et se saute d'un appui.

\begin{enumerate}
\item Pourquoi la méthode qui l'anime doit-elle rendre « vrai, je tiens
      l'écran » plutôt que de simplement dessiner ?
\item Décrivez ce qui arrive si elle ne le fait pas, dans un jeu où une IA joue.
\item Pourquoi l'ouverture doit-elle \emph{consommer} l'appui qui la saute ?
      Décrivez le comportement exact si elle ne le fait pas, du point de vue de
      l'utilisateur --- et non du code.
\end{enumerate}
\end{nklogique}
